\documentclass{amsart}
% For dealing with references we use the comment environment
\usepackage{verbatim}
\newenvironment{reference}{\comment}{\endcomment}
%\newenvironment{reference}{}{}
\newenvironment{slogan}{\comment}{\endcomment}
\newenvironment{history}{\comment}{\endcomment}

% For commutative diagrams we use Xy-pic
\usepackage[all]{xy}

% We use 2cell for 2-commutative diagrams.
\xyoption{2cell}
\UseAllTwocells

% We use multicol for the list of chapters between chapters
\usepackage{multicol}

% This is generally recommended for better output
\usepackage{lmodern}
\usepackage[T1]{fontenc}

% For cross-file-references

% Package for hypertext links:
\usepackage{hyperref}

% For any local file, say "hello.tex" you want to link to please

% Theorem environments.
%
\theoremstyle{plain}
\newtheorem{theorem}[subsection]{Theorem}
\newtheorem{proposition}[subsection]{Proposition}
\newtheorem{lemma}[subsection]{Lemma}

\theoremstyle{definition}
\newtheorem{definition}[subsection]{Definition}
\newtheorem{example}[subsection]{Example}
\newtheorem{exercise}[subsection]{Exercise}
\newtheorem{situation}[subsection]{Situation}

\theoremstyle{remark}
\newtheorem{remark}[subsection]{Remark}
\newtheorem{remarks}[subsection]{Remarks}

\numberwithin{equation}{subsection}

% Macros
%
\def\lim{\mathop{\mathrm{lim}}\nolimits}
\def\colim{\mathop{\mathrm{colim}}\nolimits}
\def\Spec{\mathop{\mathrm{Spec}}}
\def\Hom{\mathop{\mathrm{Hom}}\nolimits}
\def\Ext{\mathop{\mathrm{Ext}}\nolimits}
\def\SheafHom{\mathop{\mathcal{H}\!\mathit{om}}\nolimits}
\def\SheafExt{\mathop{\mathcal{E}\!\mathit{xt}}\nolimits}
\def\Sch{\mathit{Sch}}
\def\Mor{\mathop{\mathrm{Mor}}\nolimits}
\def\Ob{\mathop{\mathrm{Ob}}\nolimits}
\def\Sh{\mathop{\mathit{Sh}}\nolimits}
\def\NL{\mathop{N\!L}\nolimits}
\def\CH{\mathop{\mathrm{CH}}\nolimits}
\def\proetale{{pro\text{-}\acute{e}tale}}
\def\etale{{\acute{e}tale}}
\def\QCoh{\mathit{QCoh}}
\def\Ker{\mathop{\mathrm{Ker}}}
\def\Im{\mathop{\mathrm{Im}}}
\def\Coker{\mathop{\mathrm{Coker}}}
\def\Coim{\mathop{\mathrm{Coim}}}

% Boxtimes
%
\DeclareMathSymbol{\boxtimes}{\mathbin}{AMSa}{"02}

%
% Macros for moduli stacks/spaces
%
\def\QCohstack{\mathcal{QC}\!\mathit{oh}}
\def\Cohstack{\mathcal{C}\!\mathit{oh}}
\def\Spacesstack{\mathcal{S}\!\mathit{paces}}
\def\Quotfunctor{\mathrm{Quot}}
\def\Hilbfunctor{\mathrm{Hilb}}
\def\Curvesstack{\mathcal{C}\!\mathit{urves}}
\def\Polarizedstack{\mathcal{P}\!\mathit{olarized}}
\def\Complexesstack{\mathcal{C}\!\mathit{omplexes}}
% \Pic is the operator that assigns to X its picard group, usage \Pic(X)
% \Picardstack_{X/B} denotes the Picard stack of X over B
% \Picardfunctor_{X/B} denotes the Picard functor of X over B
\def\Pic{\mathop{\mathrm{Pic}}\nolimits}
\def\Picardstack{\mathcal{P}\!\mathit{ic}}
\def\Picardfunctor{\mathrm{Pic}}
\def\Deformationcategory{\mathcal{D}\!\mathit{ef}}

\setlength{\emergencystretch}{3em}
\begin{document}
\title[Stacks Project, Tag 0EX8]{475 --- Generic triviality of a line bundle extends over proper flat finitely presented valuation-ring families with integral closed fibre and one-dimensional generic global sections}
\maketitle
\noindent Copyright (C) 2005--2025 Johan de Jong. Stacks Project authors. Source: \href{https://stacks.math.columbia.edu/tag/0EX8}{Tag 0EX8}.

\noindent Editorial difficulty note (not author text). Difficulty H2: The proof identifies global sections over a valuation ring as a free rank-one lattice inside the generic sections. It then rules out vanishing of its generator on the closed fibre by a divisibility argument and uses the relative effective-Cartier criterion to show that its zero locus is empty..

\noindent Editorial provenance note (not author text). History: excerpt from revision \texttt{a0444\allowbreak 6e57e\allowbreak c1fbc\allowbreak 252a8\allowbreak 71afc\allowbreak ec775\allowbreak 2fb28\allowbreak 07b14}, more-morphisms.tex, lines 9333--9407. Reference presentation was adapted; original-author context and core support blocks were retained or added. The mathematical wording is unchanged. Licensed under GNU FDL 1.2 or later; no invariant sections or cover texts. See ../licenses/COPYING. Context blocks reproduce author definitions and supporting results; they are not additional counted problems.

\begin{lemma}
\label{lemma-triviality-generic-fibre-valuation-ring}
Let $f : X \to S$ be a proper flat morphism of finite presentation.
Let $\mathcal{L}$ be an invertible $\mathcal{O}_X$-module.
Assume
\begin{enumerate}
\item $S$ is the spectrum of a valuation ring,
\item $\mathcal{L}$ is trivial on the generic fibre $X_\eta$ of $f$,
\item the closed fibre $X_0$ of $f$ is integral,
\item $H^0(X_\eta, \mathcal{O}_{X_\eta})$ is equal to the function field of $S$.
\end{enumerate}
Then $\mathcal{L}$ is trivial.
\end{lemma}

\begin{proof}
Write $S = \Spec(A)$. We will first prove the lemma when $A$ is a
discrete valuation ring (as this is the case most often used in practice).
Let $\pi \in A$ be a uniformizer.
Take a trivializing section $s \in \Gamma(X_\eta, \mathcal{L}_\eta)$.
After replacing $s$ by $\pi^n s$ if necessary
we can assume that $s \in \Gamma(X, \mathcal{L})$.
If $s|_{X_0} = 0$, then we see
that $s$ is divisible by $\pi$ (because $X_0$ is the scheme theoretic
fibre and $X$ is flat over $A$). Thus we may assume that $s|_{X_0}$
is nonzero. Then the zero locus $Z(s)$ of $s$ is contained in $X_0$
but does not contain the generic point of $X_0$ (because $X_0$ is integral).
This means that the $Z(s)$ has codimension $\geq 2$ in $X$ which contradicts
Divisors, Lemma \href{https://stacks.math.columbia.edu/tag/0BCN}{lemma effective Cartier codimension 1 (Tag 0BCN)}
unless $Z(s) = \emptyset$ as desired.

\medskip\noindent
Proof in the general case. Since the valuation ring $A$ is
coherent (Algebra, Example \href{https://stacks.math.columbia.edu/tag/0EWV}{example valuation ring coherent (Tag 0EWV)})
we see that $H^0(X, \mathcal{L})$ is a coherent $A$-module, see
Derived Categories of Schemes, Lemma
\href{https://stacks.math.columbia.edu/tag/0EX6}{lemma cohomology over coherent ring (Tag 0EX6)}.
Equivalently, $H^0(X, \mathcal{L})$ is a finitely presented
$A$-module (Algebra, Lemma \href{https://stacks.math.columbia.edu/tag/05CX}{lemma coherent ring (Tag 05CX)}).
Since $H^0(X, \mathcal{L})$ is torsion free (by flatness of $X$ over $A$),
we see from More on Algebra, Lemma
\href{https://stacks.math.columbia.edu/tag/0ASP}{lemma generalized valuation ring modules (Tag 0ASP)}
that $H^0(X, \mathcal{L}) = A^{\oplus n}$ for some $n$.
By flat base change (Cohomology of Schemes, Lemma
\href{https://stacks.math.columbia.edu/tag/02KH}{lemma flat base change cohomology (Tag 02KH)})
we have
$$
K = H^0(X_\eta, \mathcal{O}_{X_\eta}) \cong
H^0(X_\eta, \mathcal{L}_\eta) =
H^0(X, \mathcal{L}) \otimes_A K
$$
where $K$ is the fraction field of $A$. Thus $n = 1$.
Pick a generator $s \in H^0(X, \mathcal{L})$.
Let $\mathfrak m \subset A$ be the maximal ideal.
Then $\kappa = A/\mathfrak m = \colim A/\pi$ where
this is a filtered colimit over nonzero $\pi \in \mathfrak m$
(here we use that $A$ is a valuation ring).
Thus $X_0 = \lim X \times_S \Spec(A/\pi)$.
If $s|_{X_0}$ is zero, then for some $\pi$
we see that $s$ restricts to zero on $X \times_S \Spec(A/\pi)$, see
Limits, Lemma \href{https://stacks.math.columbia.edu/tag/01Z0}{lemma descend section (Tag 01Z0)}.
But if this happens, then $\pi^{-1} s$ is
a global section of $\mathcal{L}$ which contradicts
the fact that $s$ is a generator of $H^0(X, \mathcal{L})$.
Thus $s|_{X_0}$ is not zero. Let $Z(s) \subset X$ be the zero scheme of $s$.
Since $s|_{X_0}$ is not zero and since $X_0$ is integral,
we see that $Z(s)_0 \subset X_0$ is an effective Cartier divisor.
Since $f$ is proper and $S$ is local, every point of $Z(s)$
specializes to a point of $Z(s)_0$. Thus by
Divisors, Lemma \href{https://stacks.math.columbia.edu/tag/062Y}{lemma fibre Cartier (Tag 062Y)} part (3)
we see that $Z(s)$ is a relative effective Cartier divisor,
in particular $Z(s) \to S$ is flat.
Hence if $Z(s)$ were nonemtpy, then $Z(s)_\eta$ would be nonempty
which contradicts the fact that $s|_{X_\eta}$ is a trivialization
of $\mathcal{L}_\eta$. Thus $Z(s) = \emptyset$ as desired.
\end{proof}
\end{document}
